\chapter{What We Learned}

Bootstrapping an async-first kernel on real hardware (QEMU) surfaced a series
of latent bugs that static analysis and self-tests never caught. Each fix
enabled the next layer of work. This chapter records the most instructive, in
the order they were discovered, and the design insights they produced.

\section{Bugs found on hardware}

\subsection{Panic handler recursed through heap-dependent \texttt{logln}}

\textbf{Symptom:} any fault during early boot (before the heap allocator was
ready) silently triple-faulted instead of printing a panic message --- making
\emph{every} early bug look like a mysterious hang.

\textbf{Root cause:} the panic handler used \texttt{logln!}, which routes
through \texttt{INT\_STATE} --- a \texttt{LazyLock} whose first-use
initialization allocates from the kernel heap (\texttt{alloc::vec!}). A fault
before the heap was ready caused the panic handler's \texttt{logln!} to trigger
a heap allocation $\rightarrow$ another fault $\rightarrow$ re-entered the
panic handler $\rightarrow$ infinite recursion $\rightarrow$ stack overflow
$\rightarrow$ triple fault.

\textbf{Fix:} the panic handler now uses \texttt{early\_logln!} (direct serial
writes, no heap, no \texttt{INT\_STATE}). A panic handler must never depend on
the allocator it might be reporting a failure about.

This was the single most valuable fix: it unmasked every subsequent fault as a
visible panic instead of a silent triple-fault.

\subsection{Blocked threads lost their execution context}

\textbf{Symptom:} a thread that blocked in \texttt{wait()} (or \texttt{sleep()})
restarted from its entry point when woken, instead of resuming where it
blocked. The async-syscall demo coordinator reset entirely on every wake.

\textbf{Root cause:} \texttt{block\_thread} cleared the LP scheduler's
\texttt{current\_handle} for the running thread, so the subsequent
\texttt{cond\_yield\_lp} read \texttt{curr\_tid == None} and
\texttt{switch\_ctx(null, ...)} \textbf{skipped saving} the blocking thread's
execution context.

\textbf{Fix:} a self-blocking thread now stays \texttt{current\_handle} until
the yield saves its context; \texttt{RoundRobin::next} declines to re-queue a
\texttt{Blocked} thread. This was foundational: \texttt{wait()} and
\texttt{sleep()} had never been exercised before.

\subsection{RwLock held-across-call deadlocks in \texttt{sleep()} and
\texttt{abort()}}

\textbf{Symptom:} \texttt{sleep()} worked reliably with diagnostic logging but
hung without it --- a classic heisenbug.

\textbf{Root cause:} \texttt{sleep()} held a \texttt{SYSTEM\_SCHEDULER.read()}
guard (plus the LP scheduler lock) across
\texttt{SYSTEM\_SCHEDULER.write().block\_thread(...)} --- a read-then-write
RwLock self-deadlock. The diagnostic version happened to bind the thread ID
first, releasing the guard. \texttt{abort()} had the same pattern with the LP
scheduler lock.

\textbf{Fix:} bind the value out of the scrutinee first so the temporary
guards drop before the re-locking call. This applies to \emph{any} use of
RwLock where a read guard scope is extended by an \texttt{if let} or
\texttt{match} binding.

\subsection{Quantum timer grew the queue without bound}

\textbf{Symptom:} after extended operation, the per-LP timer queue contained
thousands of quantum events, consuming unbounded memory.

\textbf{Root cause:} \texttt{clear\_ctx\_switch\_pending} enqueued a new
quantum \texttt{TimerEvent} on every yield, not just on quantum expiry. Manual
yields accumulated quantum events forever.

\textbf{Fix:} an \texttt{armed} flag keeps exactly one quantum event in flight
(the quantum's own firing clears it). Also, \texttt{cond\_yield\_lp} re-arms
even when there is nothing to switch to (previously the timer stopped when only
one thread was runnable), and \texttt{process\_events} stops the
level-triggered timer when the queue drains.

\subsection{A thread cannot free its own kernel stack}

\textbf{Symptom:} a returning kernel thread hung during teardown.

\textbf{Root cause:} \texttt{ThreadContext::Drop} deallocates the kernel stack
--- but the returning thread is still executing on that stack when
\texttt{abort} drops it. The CPU eventually faults on freed memory.

\textbf{Fix:} a \textbf{deferred reaper}. \texttt{abort} moves the dying thread
to \texttt{DEAD\_THREADS} (\texttt{IdTable::take\_element} extracts it without
dropping), and \texttt{cond\_yield\_lp} calls \texttt{reap\_dead\_threads()}
\emph{after} switching away, so the stack is freed from a different thread's
context.

\subsection{Stack allocator never registered guard pages}

\textbf{Symptom:} \texttt{deallocate\_stack} always failed with
\texttt{InvalidStack}, causing the reaper to leak every freed stack.

\textbf{Root cause:} \texttt{allocate\_stack} never populated
\texttt{KERNEL\_GUARD\_PAGE\_SET}, so \texttt{validate\_stack} could never find
the guards.

\textbf{Fix:} register lower and upper guard pages on allocation and rewrite
\texttt{deallocate\_stack} to derive the page count from them. Reference-count
guard pages (\texttt{BTreeMap<VAddr, usize>}) so adjacent stacks sharing a
guard survive until both are freed.

\subsection{LA57 paging-mode mismatch on x86-64}

\textbf{Symptom:} a \texttt{\#GP} with \texttt{RAX = 0xff90000000000000}
(non-canonical under 4-level paging) during kernel-heap large-page mapping.

\textbf{Root cause:} \texttt{get\_vaddr\_sig\_bits()} read the CPU's
\emph{maximum} linear address width from CPUID. With \texttt{-cpu max} that is
57, so the kernel selected the 57-bit address map. But Limine boots with
4-level paging (48-bit), where addresses in the 57-bit map are non-canonical.

\textbf{Fix:} derive the address width from \texttt{CR4.LA57} (the
\emph{active} paging mode), not from CPUID capability.

\subsection{ASID-0 collision}

\textbf{Symptom:} the first user thread faulted at its entry point with
an instruction abort (\texttt{EC=0x21}, ``no change in exception level''),
despite correct page table entries.

\textbf{Root cause:} the global address-space table started empty, so the first
user AS got id 0, which equals \texttt{KERNEL\_ASID}. \texttt{Thread::new}
built a kernel thread (EL1) instead of a user thread (EL0), running the user
code at EL1 where \texttt{PXN=1} forbids execution. The EC decode
(\texttt{0x21} = fault without EL change) gave the diagnosis: the code was
never at EL0.

\textbf{Fix:} pre-populate the table with the kernel AS at id 0, so user
spaces always get non-zero ids. A second bug in the same test:
\texttt{sync\_dispatcher} added 4 to \texttt{ELR\_EL1} on SVC, but AArch64
already sets ELR to the instruction \emph{after} SVC --- the +4 skipped the
user's loop instruction into zeroed memory.

\subsection{x86-64 trampoline halted on thread return}

\textbf{Symptom:} on x86-64, the async demo's worker completed but the
coordinator never resumed.

\textbf{Root cause:} the x86-64 \texttt{kernel\_thread\_trampoline} entered a
\texttt{hlt} loop when a thread's entry returned, instead of calling
\texttt{abort()} like the AArch64 trampoline. The CPU halted, freezing the LP.

\textbf{Fix:} make the x86-64 trampoline call \texttt{abort} on entry return.

\section{Key insights}

\begin{enumerate}
    \item \textbf{The block-yield-wake cycle is the hard problem.} Designing
        the completion ABI was straightforward compared to making
        \texttt{wait()} correctly save and restore thread state through the
        cooperative scheduler. Every scheduling primitive (\texttt{sleep},
        \texttt{wait}, \texttt{abort}) exercised a new path and surfaced a new
        bug. The block-yield-wake path was untested before the async demo.

    \item \textbf{Booting on hardware finds bugs that static analysis cannot.}
        The panic-handler recursion, ASID-0 collision, and LA57 paging
        mismatch were invisible in self-tests and logical reasoning. Every one
        was found by booting QEMU and observing the actual fault register
        values.

    \item \textbf{The observer/waker pattern is the right primitive.} Once
        \texttt{Completion} was made \texttt{Observable} and \texttt{wait()}
        used the same path as \texttt{sleep()}, the entire completion subsystem
        fell into place. The exit-observer mechanism is the natural endpoint:
        a thread exiting \emph{is} the completion event. The worker never
        references the capability.

    \item \textbf{The async model genuinely removes the retrofit tax.} On
        CharlotteOS, \texttt{wait(cap)} blocks exactly as \texttt{sleep(50ms)}
        blocks --- both register a \texttt{Waker} on an \texttt{Observable} and
        yield. There is no emulation, no thread pool, no signal-handler
        tightrope. The kernel and the runtime agree about what is fundamental.

    \item \textbf{The sitas traits were essential for co-design.} The
        \texttt{ReactorBackend} trait became the executable specification for
        the async syscall ABI. The sitas-charlotte backend validates the ABI
        shapes before any kernel path exists. This bidirectional validation is
        the core of the collaboration.

    \item \textbf{The deferred reaper is necessary for any thread that exits.}
        A thread cannot free its own stack. The reaper pattern (move dead
        threads to a zombie list, drop them from another context) is the
        minimum safe discipline.

    \item \textbf{\texttt{RwLock} held-across-call is a systemic hazard.} The
        Rust pattern of \texttt{if let Some(val) = lock().method()} extends the
        guard's lifetime across the body, which is invisible to the reader.
        Bind the value first, then use it. This caused both the sleep heisenbug
        and the abort deadlock.
\end{enumerate}

\endinput
